\documentclass[12pt]{article}
\usepackage{amsmath, amssymb}
\usepackage[margin=1in]{geometry}
\setlength{\parskip}{6pt}
\title{QUBO Formulation: Top Eigencentrality Node Selection Problem}
\date{}
\begin{document}
\maketitle

\subsection*{Top Eigencentrality Node Selection Problem}

\paragraph{Problem Statement.}
Given a graph $G=(V, E)$ with $N = |V|$ nodes, a target integer $\tau \in \{1, \dots, N\}$, and pre-computed eigencentrality scores $c_i$ for each node $i \in V$, the Top Eigencentrality Node Selection Problem seeks to identify a subset of exactly $\tau$ nodes that maximizes the sum of their eigencentrality scores. Formally, the problem is to find a subset $S \subseteq V$ with cardinality $|S| = \tau$ such that the objective function $\sum_{i \in S} c_i$ is maximized.

\paragraph{Decision Variables.}
We introduce binary decision variables $x_i \in \{0, 1\}$ for each node $i \in \{0, 1, \dots, N-1\}$. The variable $x_i$ takes the value $1$ if node $i$ is selected in the solution, and $0$ otherwise.

\paragraph{QUBO Derivation.}
We convert the maximization problem into a minimization problem suitable for QUBO formulation by negating the objective. The original problem is equivalent to:
\begin{align}
    \min_{x \in \{0,1\}^N} \left( -\sum_{i=0}^{N-1} c_i x_i \right)
    \quad \text{subject to} \quad
    \sum_{i=0}^{N-1} x_i = \tau.
\end{align}
We incorporate the equality constraint into the objective using a quadratic penalty term with weight $\lambda > 0$. The penalized Hamiltonian $H(x)$ is given by:
\begin{align}
    H(x) = -\sum_{i=0}^{N-1} c_i x_i + \lambda \left( \sum_{i=0}^{N-1} x_i - \tau \right)^2.
\end{align}
Expanding the squared penalty term yields:
\begin{align}
    \lambda \left( \sum_{i=0}^{N-1} x_i - \tau \right)^2
    &= \lambda \left( \sum_{i=0}^{N-1} x_i^2 + \sum_{i \neq j} x_i x_j - 2\tau \sum_{i=0}^{N-1} x_i + \tau^2 \right).
\end{align}
Using the idempotent property of binary variables, $x_i^2 = x_i$, we simplify the expression:
\begin{align}
    \lambda \left( \sum_{i=0}^{N-1} x_i + \sum_{i \neq j} x_i x_j - 2\tau \sum_{i=0}^{N-1} x_i + \tau^2 \right)
    &= \lambda \sum_{i=0}^{N-1} (1 - 2\tau) x_i + \lambda \sum_{i \neq j} x_i x_j + \lambda \tau^2.
\end{align}
Noting that $\sum_{i \neq j} x_i x_j = 2 \sum_{0 \leq i < j \leq N-1} x_i x_j$, the penalty contribution becomes:
\begin{align}
    \lambda \sum_{i=0}^{N-1} (1 - 2\tau) x_i + 2\lambda \sum_{0 \leq i < j \leq N-1} x_i x_j + \lambda \tau^2.
\end{align}
Combining the objective and penalty terms, the full QUBO Hamiltonian is:
\begin{align}
    H(x) = \sum_{i=0}^{N-1} \left( \lambda(1 - 2\tau) - c_i \right) x_i + 2\lambda \sum_{0 \leq i < j \leq N-1} x_i x_j + \lambda \tau^2.
\end{align}
Reading off the coefficients from the general form $H(x) = \sum_{i} Q_{i,i} x_i + \sum_{i < j} Q_{i,j} x_i x_j + c$, we obtain the QUBO parameters:
\begin{align}
    Q_{i,i} &= \lambda(1 - 2\tau) - c_i, \quad \forall i \in \{0, \dots, N-1\}, \\
    Q_{i,j} &= 2\lambda, \quad \forall i, j \in \{0, \dots, N-1\}, \; i < j, \\
    c &= \lambda \tau^2.
\end{align}

\paragraph{Penalty Weight Justification.}
The penalty weight $\lambda$ must be chosen sufficiently large to ensure that the global minimum of $H(x)$ corresponds to a feasible solution. Specifically, $\lambda$ must exceed the maximum possible gain in the objective function achievable by violating the constraint. Since the objective coefficients are $c_i$, the maximum gain from selecting an additional node (or deviating from the count $\tau$) is bounded by $\max_{i} c_i$. Thus, setting $\lambda > \max_{i \in \{0, \dots, N-1\}} c_i$ guarantees that any infeasible solution incurs a penalty strictly greater than any potential improvement in the objective value, forcing the optimizer to satisfy $\sum_{i} x_i = \tau$.

\paragraph{Correctness Argument.}
Minimizing $H(x)$ over $x \in \{0,1\}^N$ recovers the optimal solution to the original problem. For any feasible assignment where $\sum_{i} x_i = \tau$, the penalty term vanishes, and $H(x)$ reduces to the negative objective function. For any infeasible assignment, the penalty term is strictly positive. By the choice of $\lambda$, the energy of any infeasible solution exceeds that of the best feasible solution. Consequently, the global minimum of $H(x)$ must occur at a feasible point that minimizes the negative objective, which is equivalent to maximizing the original objective.

\paragraph{Illustrative Example.}
Consider an instance with $N=3$ nodes, target $\tau=1$, and eigencentrality scores $c_0=5, c_1=3, c_2=4$. We choose $\lambda=6$, satisfying $\lambda > \max(c_i)=5$. Substituting these values into the general formulas:
\begin{align}
    Q_{0,0} &= 6(1 - 2) - 5 = -11, \\
    Q_{1,1} &= 6(1 - 2) - 3 = -9, \\
    Q_{2,2} &= 6(1 - 2) - 4 = -10, \\
    Q_{0,1} &= Q_{0,2} = Q_{1,2} = 2(6) = 12, \\
    c &= 6(1)^2 = 6.
\end{align}
The resulting Hamiltonian is:
\begin{align}
    H(x) = -11 x_0 - 9 x_1 - 10 x_2 + 12 x_0 x_1 + 12 x_0 x_2 + 12 x_1 x_2 + 6.
\end{align}
The feasible solutions are the permutations of $(1,0,0)$. Evaluating $H(x)$:
\begin{align}
    H(1,0,0) &= -11 + 6 = -5, \\
    H(0,1,0) &= -9 + 6 = -3, \\
    H(0,0,1) &= -10 + 6 = -4.
\end{align}
The minimum energy is $-5$ at $x=(1,0,0)$, corresponding to selecting node $0$ with the maximum score $5$.
\end{document}